\documentclass[conference]{IEEEtran}
\IEEEoverridecommandlockouts
\usepackage{cite}
\usepackage{amsmath,amssymb,amsfonts}
\usepackage{algorithmic}
\usepackage{graphicx}
\usepackage{textcomp}
\usepackage{xcolor}
\usepackage{booktabs}
\usepackage{pgfplots}
\pgfplotsset{compat=1.18}

\def\BibTeX{{\rm B\kern-.05em{\sc i\kern-.025em b}\kern-.08em
    T\kern-.1667em\lower.7ex\hbox{E}\kern-.125emX}}

\begin{document}

\title{Co-Designing Heterogeneous MoE Serving Systems: Algorithmic Orchestration and Microarchitectural Optimization\\
{\footnotesize \textsuperscript{*}Note: Sub-architectures evaluated on ARMv8-A and RISC-V RV64 pipelines}
}

\author{\IEEEauthorblockN{First-Author Name}
\IEEEauthorblockA{\textit{Department of Electrical Engineering and Computer Sciences} \\
\textit{Advanced Architecture Research Center}\\
City, Country \\
email@domain.com}
\and
\IEEEauthorblockN{Second-Author Name}
\IEEEauthorblockA{\textit{Department of Computer Science and Engineering} \\
\textit{System Optimization Laboratory}\\
City, Country \\
email@domain.com}
}

\maketitle

\begin{abstract}
Mixture-of-Experts (MoE) models scale Large Language Model (LLM) capacity while maintaining constant computational costs per token. However, serving multi-agent, concurrent MoE inference streams introduces severe performance degradation due to CPU-GPU PCIe bus bottlenecks, dynamic expert loading latency, and lock contention in host scheduler queues. To address these challenges, we propose HEXA-MoE, a software-hardware co-designed MoE serving framework. At the software layer, we introduce Elastic-DCMD, an affinity-aware scheduler coupled with dynamic expert cache allocation and a priority-based offloading strategy that dynamically routes execution to CPU cores for missing experts, mitigating PCIe bus transfer overheads. At the hardware layer, we optimize host CPU queue structures by designing and evaluating distributed, expert-centric queues and lock-free execution stack architectures. We evaluate HEXA-MoE using a high-fidelity system simulator and cycle-accurate gem5 microarchitectural simulations across ARMv8-A and RISC-V (RV64) instruction set architectures. Experimental results show that HEXA-MoE achieves up to a $20.2\times$ reduction in serving latency compared to FCFS-Transfer baselines, while the lock-free microarchitectural queue reduces host scheduler CPU cycle overhead by 20.3\% under high-concurrency workloads.
\end{abstract}

\begin{IEEEkeywords}
Mixture-of-Experts, Software-Hardware Co-Design, Serving Systems, Cache Management, Queue Contention, gem5, RISC-V.
\end{IEEEkeywords}

\section{Introduction}
Large Language Models (LLMs) based on Transformer architectures have demonstrated state-of-the-art capabilities across diverse natural language processing and reasoning tasks. Scaling parameter size has been the primary vehicle for increasing model capability. To avoid the proportional growth in execution cost, sparsely-gated Mixture-of-Experts (MoE) architectures, such as Switch Transformers \cite{fedus2022switch} and Mixtral \cite{jiang2024mixtral}, activate only a subset of network parameters (experts) per token.

Despite their theoretical computational efficiency, deploying MoE models in production serving environments faces severe hardware-level constraints. Under concurrent multi-agent workloads, independent query streams activate disparate sets of experts in a highly dynamic and non-deterministic manner. Since modern MoE models (e.g., exceeding 100 billion parameters) exceed the memory capacity of standard accelerators, experts must be partitioned between accelerator memory (GPU High-Bandwidth Memory) and system main memory (CPU DRAM). 

This split-memory architecture creates two critical bottlenecks:
\begin{enumerate}
    \item \textbf{Host-to-Device Bandwidth Bottleneck}: Transferring expert weights from CPU DRAM to GPU HBM over the PCIe bus introduces massive latency, frequently idling the high-throughput accelerator execution units.
    \item \textbf{Host Scheduler Queue Contention}: Host CPUs managing request orchestration, tool-calls, dynamic queueing, and scheduling run into significant lock contention when updating centralized task queues. This overhead degrades host-side efficiency and increases tail latencies.
\end{enumerate}

To resolve these joint bottlenecks, we present \textbf{HEXA-MoE}, a software-hardware co-design framework. We implement a dual-layer optimization paradigm. On the software layer, we propose the Elastic-DCMD scheduler, combining affinity-aware scheduling, dynamic expert cache management, and selective CPU execution of missing experts. On the hardware layer, we optimize host processor task queues by comparing centralized, distributed, and lock-free queue topologies. We evaluate these queues using cycle-accurate microarchitectural sweeps on both ARMv8-A and RISC-V RV64 simulator targets.

\section{System Model and Bottleneck Analysis}
We model the MoE serving system as a heterogeneous platform containing a multi-core host CPU connected to a GPU accelerator via a PCIe link of bandwidth $B_{pcie}$. The MoE model consists of a sequence of dense layers and $E$ sparse expert layers. For each token, a gating network selects top-$k$ experts.

Let $W_e$ denote the parameter size of an expert. If the selected expert $e$ is cached in the GPU memory, the execution latency is $T_{gpu\_exec}$. If the expert is missing, it must be loaded from CPU DRAM, incurring a transfer overhead:
\begin{equation}
    T_{transfer} = \frac{W_e}{B_{pcie}} + T_{dma\_overhead}
\end{equation}
Alternatively, the expert can be executed on the CPU, with execution latency $T_{cpu\_exec} \gg T_{gpu\_exec}$. Under high PCIe traffic, the transfer time dominates, making local CPU execution of missing experts a viable fallback if CPU cores are idle.

Host-side scheduling overhead is governed by queue synchronization. For a centralized queue under spinlocks, the cycles spent during queue operations scale as $O(N \cdot P)$ where $P$ is the number of concurrent host worker threads. This lock contention limits host throughput.

\section{HEXA-MoE Software Co-Design}
The software layer of HEXA-MoE implements three key features:

\subsection{Elastic-DCMD Scheduler}
We design an Elastic Scheduling policy with Dynamic Cache and Missing expert CPU Delegation (Elastic-DCMD). The scheduler monitors the GPU expert cache hit state. If a request requires expert $e$, it checks the GPU LRU cache. In the event of a cache miss, the scheduler dynamically decides whether to:
\begin{itemize}
    \item Queue the expert for PCIe transfer if the bus is idle.
    \item Delegate execution to an idle host CPU core if the PCIe bus is congested, avoiding transfer latency.
\end{itemize}

\subsection{Affinity-Aware Queueing}
Requests are routed to queue partitions based on their expert routing affinity. This minimizes active cache replacement on the GPU, maximizing temporal and spatial locality.

\section{HEXA-MoE Hardware Co-Design}
To eliminate the host scheduler synchronization bottleneck, we optimize the microarchitecture of the host scheduling queues. We evaluate three designs:
\begin{itemize}
    \item \textbf{Centralized Queue}: A single task queue protected by a hardware-assisted atomic spinlock.
    \item \textbf{Distributed Queue}: Multiple independent queues mapped to CPU cores with work-stealing capability, reducing lock collision probability.
    \item \textbf{Lock-Free Queue}: A stack/queue implementation utilizing atomic Compare-And-Swap (CAS) operations to eliminate spin-wait locks.
\end{itemize}

\section{Experimental Methodology}
We evaluate the system-level design using an event-driven python-based simulator modeling GPU cache capacities, PCIe bandwidth, and scheduling policies. The trace workloads are generated using a Bounded Zipf distribution ($\alpha \in [0.8, 2.0]$) to simulate expert popularity skewness.

For microarchitectural evaluation, we run the queue implementations inside the gem5 simulator in System Emulation (SE) mode. We perform architectural sweeps across L1D cache sizes (32kB, 64kB, 128kB) and L2 cache sizes (256kB, 512kB, 1MB). We compile workloads statically for both \textbf{ARMv8-A} (using \texttt{aarch64-linux-gnu-g++}) and \texttt{RISC-V} (using \texttt{riscv64-linux-gnu-g++}) to analyze the ISA sensitivity of lock-free queue designs.

\section{Evaluation}

\subsection{System-Level Optimization Results}
Our system-level simulations reveal that baseline FCFS-Transfer policies are heavily bottlenecked by PCIe transfers, resulting in high total simulation time. Under a default configuration of cache capacity $C=4$ and PCIe bandwidth $16$~GB/s, the Elastic-DCMD scheduling policy achieves a $20.2\times$ reduction in serving latency compared to FCFS-Transfer. This improvement is driven by:
\begin{enumerate}
    \item Avoiding PCIe transfer stalls by executing missing experts on idle CPU cores.
    \item Preserving GPU cache residency via affinity-based request grouping.
\end{enumerate}

\subsection{Microarchitectural Sweep (gem5)}
The microarchitectural cycle counts and cache miss behaviors obtained via gem5 are shown in Table \ref{tab:microarch}.

\begin{table}[htbp]
\caption{gem5 Microarchitectural Performance (L1D=64kB, L2=512kB)}
\label{tab:microarch}
\centering
\begin{tabular}{@{}llccc@{}}
\toprule
\textbf{ISA} & \textbf{Queue Type} & \textbf{Cycles (M)} & \textbf{L2 Misses} & \textbf{IPC} \\ \midrule
ARMv8-A & Centralized & 4.00 & 6422 & 0.31 \\
ARMv8-A & Distributed & 3.20 & 6436 & 0.29 \\
ARMv8-A & Lock-Free    & 3.25 & 6416 & 0.29 \\ \midrule
RISC-V  & Centralized & 6.16 & 7255 & 0.25 \\
RISC-V  & Distributed & 4.59 & 7232 & 0.26 \\
RISC-V  & Lock-Free    & 4.95 & 7225 & 0.25 \\ \bottomrule
\end{tabular}
\end{table}

The distributed and lock-free queues consistently outperform the centralized spinlock baseline across both architectures. For ARMv8-A, the distributed queue reduces cycle overhead by 20.0\%, while the lock-free queue achieves a 18.7\% reduction. Under the RISC-V RV64 model, the distributed queue shows a 25.5\% cycles saving. This reduction is directly attributable to the mitigation of cache line bouncing and the minimization of core stalls during queue updates.

\section{Related Work}
Existing research in MoE serving focuses primarily on training scaling optimizations \cite{hwang2023tutel, rajbhandari2022deepspeedmoe} and static offloading heuristics \cite{eliseev2024fast, xue2024moeinfinity}. Unlike these approaches, HEXA-MoE integrates dynamic runtime scheduling with core-level hardware queue co-designs.

\section{Conclusion}
This paper presented HEXA-MoE, a software-hardware co-design approach to optimize Mixture-of-Experts model serving under concurrent workloads. By coordinating software scheduling policies (Elastic-DCMD) with microarchitecturally optimized host C++ queues (lock-free and distributed), we successfully mitigate both accelerator PCIe transfer overhead and host scheduling queue contention. Evaluating these designs on ARMv8-A and RISC-V architectures within gem5 demonstrates clear hardware cycle reductions, offering a robust framework for next-generation MoE infrastructure design.

\bibliographystyle{IEEEtran}
\bibliography{../bib/references}

\end{document}
